\documentclass[10pt,a4paper]{amsart}
\usepackage[margin=19mm]{geometry}
\usepackage{amsmath,amssymb,mathtools}
\usepackage[scr=rsfs,cal=euler]{mathalfa}
\usepackage{xfrac}
\usepackage[unicode,hidelinks,bookmarks=false]{hyperref}
\newcommand{\ie}{i.e.\ }
\newcommand{\cf}{cf.\ }
\newcommand{\resp}{respectively\ }
\newcommand{\Subsection}{\subsection}
\providecommand{\nobreakdash}{\nobreak\hbox{-}}
\setlength{\emergencystretch}{2em}
\multlinegap0pt
\usepackage[shortlabels]{enumitem}
\usepackage{amssymb}
\newtheorem{lemma}{Lemma}
\newcommand{\bbe}{\mathbb{E}}
\providecommand{\diff}{\mathop{}\!\mathrm{d}}
\newcommand{\eps}{\varepsilon}
\title{Existence of a classical solution to the integro-differential equation arising in the Cram\'er--Lundberg non-life insurance model with proportional investment}
\author{Platon Promyslov}
\date{Source: arXiv:2604.05143v1 (CC BY 4.0)}
\begin{document}
\maketitle

\noindent\emph{Source and licence.} Existence of a classical solution to the integro-differential equation arising in the Cram\'er--Lundberg non-life insurance model with proportional investment. Version arXiv:2604.05143v1 (CC BY 4.0), distributed by arXiv under the Creative Commons Attribution 4.0 International licence, http://creativecommons.org/licenses/by/4.0/, as recorded in the arXiv licence metadata for this exact version and shown on the abstract page https://arxiv.org/abs/2604.05143v1. Full licence notice: \texttt{licenses/arxiv-2604.05143-CC-BY-4.0.txt}.\par\smallskip

\noindent\textit{Editorial scope.} Counted unit: the article's lemma lem:iis (source0.tex lines 251--253). The article's complete original proof of it is delivered with it (source0.tex lines 254--299, one \texttt{proof} environment of 46 lines). No mathematics is written, repaired or re-derived: the statement and the proof are the article's own lines, in the article's own notation, and no retained line is substituted or re-flowed.  Retained uncounted as the source-local context the proof needs: lines 76--81; lines 138--140; lines 204--206.  Nothing was omitted from any retained span, so the package carries no omission record.  The retained lines cite 1 works, whose entries are transcribed verbatim from the article's own bibliography source in the e-print and delivered with short editorial labels recorded in the item's reference mapping.  No retained line contains a figure environment, an image-inclusion command or drawing code, so no image is delivered and the package declares no asset dependency.  Editorial formatting only, in addition: an AMS \texttt{amsart} preamble that restates the article's own declarations and macros used by the retained lines, namely the environments \texttt{lemma} on separate counters, together with the standard packages those lines need.  The retained environments print in this document as Lemma 1 in order of appearance, whereas the article prints the same items with the numbering of its own full document.  The complete native source of the article as distributed in the arXiv e-print for this exact version is bundled unchanged as \texttt{sources/197931/source0.tex}.
The analysis relies on the following basic assumptions:
\begin{enumerate}[label=(\textbf{A\arabic*})]
    \item \label{assump:1} The distribution function $F$ is continuous at zero: $F(0)=0$.
    \item \label{assump:2} There exists a constant $\eps > 0$ such that $\bbe[\xi^\eps] < \infty$.
    \item \label{assump:3} The inequality $\gamma > 1$ holds; otherwise $\Psi(u) = 1$ (see \cite{Paulsen1993}[Theorem 3.1]).
\end{enumerate}

\begin{equation}\label{eq:volterra_zero}
g(u) = \frac{\mu}{u^\gamma e^{-\alpha/u}} \int_0^u t^{\gamma-2} e^{-\alpha/t} \big( \Phi(0+) \bar{F}(t) + (Bg)(t) \big) \diff t, \quad u > 0.
\end{equation}

\begin{lemma}[Uniform Boundedness]\label{lem:boundedness}
Suppose condition \ref{assump:2} holds. Then the continuous solution $g(u)$ is uniformly bounded on $[0,\infty)$.
\end{lemma}

\begin{lemma}[Absolute Integrability]\label{lem:iis}
Suppose conditions \ref{assump:2} and \ref{assump:3} hold, with $\varepsilon \in (0,1)$ chosen such that $\gamma-1-\varepsilon > 0$. Then $g(u) = O(u^{-1-\varepsilon})$ as $u \to \infty$. In particular, $g \in L^1(\mathbb{R}_+)$.
\end{lemma}

\begin{proof}
From Lemma \ref{lem:boundedness}, we know that $|g(u)| \le C_0$ for all $u\ge0$. Assume that for some $\delta \in [0,1]$ we have
\[
|g(u)| \le C (1+u)^{-\delta} \qquad \text{for all } u \ge 0.
\]
For $\delta = 0$ this holds with $C = C_0$. We show that there then exists a constant $C'$, independent of $u$, such that $|g(u)| \le C' (1+u)^{-(\delta+\varepsilon)}$ for all sufficiently large $u$.

We estimate the convolution $(Bg)(t)$ for large $t$. We split the integral:
\[
|(Bg)(t)| \le \int_0^{t/2} |g(t-y)| \bar F(y) \diff y + \int_{t/2}^t |g(y)| \bar F(t-y) \diff y.
\]
For the first integral, $y \le t/2$, so $t-y \ge t/2$ and $|g(t-y)| \le C (1+t/2)^{-\delta} \le C_1 t^{-\delta}$. Given that $\int_0^{t/2} \bar F(y) \diff y \le J(t) \le C_2 t^{1-\varepsilon}$ (see the proof of Lemma \ref{lem:boundedness}), the first integral does not exceed $C_3 t^{1-\delta-\varepsilon}$. For the second integral: for $y \ge t/2$ we have $|g(y)| \le C_1 t^{-\delta}$, and $\int_{t/2}^t \bar F(t-y) \diff y = \int_0^{t/2} \bar F(z) \diff z \le C_2 t^{1-\varepsilon}$. Consequently, the second integral is also bounded by $C_3 t^{1-\delta-\varepsilon}$. Thus, there exists a constant $K_1$ such that for all sufficiently large $t$ we have
\[
|(Bg)(t)| \le K_1 t^{1-\delta-\varepsilon}.
\]

Moreover, for $t \ge 1$ and since $1-\delta \ge 0$, we have $\bar F(t) \le C_0 t^{-\varepsilon} \le C_0 t^{1-\delta-\varepsilon}$. Substituting these estimates into \eqref{eq:volterra_zero}, we obtain for large $u$:
\begin{multline*}
|g(u)| \le \frac{\mu}{u^\gamma e^{-\alpha/u}} \int_0^u t^{\gamma-2} e^{-\alpha/t} \bigl( \Phi(0+)\bar F(t) + |(Bg)(t)| \bigr) \diff t \le
\\
\le \frac{\mu K_2}{u^\gamma e^{-\alpha/u}} \int_0^u t^{\gamma-2} e^{-\alpha/t} t^{1-\delta-\varepsilon} \diff t
= \frac{\mu K_2}{u^\gamma e^{-\alpha/u}} \int_0^u t^{\gamma-1-\delta-\varepsilon} e^{-\alpha/t} \diff t.
\end{multline*}
The asymptotics of the last integral as $u\to\infty$ is given by:
\[
\int_0^u t^{\gamma-1-\delta-\varepsilon} e^{-\alpha/t} \diff t \sim \frac{u^{\gamma-\delta-\varepsilon}}{\gamma-\delta-\varepsilon} e^{-\alpha/u},
\]
which is easily derived using L'Hôpital's rule. Taking into account the cancellation of factors, we get:
\[
|g(u)| \le \frac{\mu K_2}{\gamma-\delta-\varepsilon} u^{-(\delta+\varepsilon)} + o(u^{-(\delta+\varepsilon)}).
\]
Hence, for all sufficiently large $u$ we have $|g(u)| \le C' u^{-(\delta+\varepsilon)}$, as required.

Starting from $\delta = 0$, after $N = \lceil 1/\varepsilon \rceil$ steps we reach $\delta \ge 1$. We apply one more step, starting from $\delta = 1$. In this case $|g(u)| \le C u^{-1}$ and, as shown above, $|(Bg)(t)| \le K_3 t^{-\varepsilon}$. Then
\begin{multline*}
|g(u)| \le \frac{\mu}{u^\gamma e^{-\alpha/u}} \int_0^u t^{\gamma-2} e^{-\alpha/t} \bigl( \Phi(0+)\bar F(t) + K_3 t^{-\varepsilon} \bigr) \diff t
\le
\\
\le \frac{\mu K_4}{u^\gamma e^{-\alpha/u}} \int_0^u t^{\gamma-2-\varepsilon} e^{-\alpha/t} \diff t.
\end{multline*}
The asymptotics yields:
\[
|g(u)| \le \frac{\mu K_4}{\gamma-1-\varepsilon} u^{-1-\varepsilon} + o(u^{-1-\varepsilon})
\]
for all sufficiently large $u$. Thus, $g(u) = O(u^{-1-\varepsilon})$, which, since $\varepsilon>0$, implies $g \in L^1(\mathbb{R}_+)$.
\end{proof}

\begin{thebibliography}{99}
\bibitem[Paulse]{Paulsen1993} Paulsen, J.: Risk theory in a stochastic economic environment. Stochastic Process. Appl. \textbf{46}(2), 327--361 (1993)
\end{thebibliography}
\end{document}
